\documentclass{ar-1col}
\usepackage{lmodern}
\usepackage[T1]{fontenc}
\usepackage[a4paper, top=2.2cm, bottom=2.2cm, left=2.5cm, right=2.5cm]{geometry}
\usepackage{url}
\usepackage[numbers]{natbib}
\usepackage{indentfirst}
\usepackage{xcolor}
\usepackage{amsmath}
\usepackage{amssymb}
\usepackage{mathtools}
% \usepackage{physics}  % not installed; not needed
% \usepackage{bm}
% \usepackage{siunitx}
\usepackage{graphicx}
\usepackage{float}
\usepackage{needspace}
\usepackage{subcaption}
\usepackage{booktabs}
\usepackage{multirow}
\usepackage{hyperref}
% \usepackage{CJKutf8}  % removed: not installed; report is English-only
\linespread{1.2}
\setlength{\floatsep}{6pt plus 2pt minus 2pt}
\setlength{\textfloatsep}{8pt plus 2pt minus 4pt}
\setlength{\intextsep}{6pt plus 2pt minus 2pt}
\setlength{\abovecaptionskip}{4pt}
\setlength{\belowcaptionskip}{2pt}

\setcounter{secnumdepth}{4}

\jname{Research Report --- TTS Voice Cloning}
\jvol{Mar.2026}
\jyear{2026}

\begin{document}

\title{Full-Pipeline Chinese TTS Voice Cloning: \\ From Data Collection to Accelerated Deployment}

\author{Author Yunlin He$^1$
\affil{$^1$AI \& Entrepreneurship, Hong Kong University of Science and Technology, email: yuzuhe2023@163.com}}

\begin{abstract}
This report presents a complete pipeline for building a high-quality Chinese text-to-speech (TTS) voice cloning system targeting a specific speaker (``Qinche''). We cover the full lifecycle: audio data collection from online sources, multi-stage preprocessing (source separation, speaker diarization, VAD segmentation, ASR transcription), supervised fine-tuning (SFT) of Qwen3-TTS-1.7B, and comprehensive evaluation against Fish Audio S2 Pro in zero-shot mode. We investigate multiple inference acceleration strategies including CUDA Graph capture (FasterQwen3TTS), \texttt{torch.compile}, INT8 quantization, and speculative decoding. Our best configuration --- Qwen3-TTS SFT v5 epoch 3 with FasterQwen3TTS --- achieves a speaker similarity (SIM\_gt) of 0.689, word error rate (WER) of 4.25\%, and real-time factor (RTF) of 0.357 on a single A100-80GB GPU, representing a \textbf{7.1$\times$ speedup} over native inference with negligible quality loss.
\end{abstract}

\begin{keywords}
Text-to-Speech, Voice Cloning, Qwen3-TTS, Fish Audio S2 Pro, CUDA Graph, Fine-tuning, Inference Acceleration
\end{keywords}

\maketitle
\tableofcontents

%% ============================================================
\section{INTRODUCTION}
%% ============================================================

Text-to-speech (TTS) voice cloning aims to synthesize speech that faithfully reproduces a target speaker's vocal characteristics. Recent advances in large-scale neural codec language models --- such as Qwen3-TTS, Fish Audio S2, VALL-E, and CosyVoice --- have dramatically improved synthesis quality, enabling few-shot and even zero-shot voice cloning from short reference audio clips.

This report documents a complete end-to-end project that:
\begin{enumerate}
    \item Collects and processes raw audio data from online sources (Bilibili) through a multi-stage pipeline;
    \item Fine-tunes the Qwen3-TTS-1.7B model with supervised fine-tuning (SFT) on $\sim$28 minutes of curated target speaker audio;
    \item Evaluates the fine-tuned model against Fish Audio S2 Pro (5B parameters, zero-shot only);
    \item Explores and benchmarks multiple inference acceleration strategies, ultimately achieving 7.1$\times$ real-time speedup via CUDA Graph capture.
\end{enumerate}

\begin{issues}[REPORT STRUCTURE]
\begin{enumerate}
    \item \textbf{Section 2}: Model architectures and design choices
    \item \textbf{Section 3}: Data collection and preprocessing pipeline
    \item \textbf{Section 4}: Training methodology and experiment iterations
    \item \textbf{Section 5}: Evaluation framework and comprehensive results
    \item \textbf{Section 6}: Inference acceleration strategies
    \item \textbf{Section 7}: Conclusions and future work
\end{enumerate}
\end{issues}


%% ============================================================
\section{MODEL ARCHITECTURES}
%% ============================================================

\subsection{Qwen3-TTS (Primary Model)}

Qwen3-TTS is a Transformer decoder-based TTS model with a 16-codebook Residual Vector Quantization (RVQ) codec operating at 12 Hz frame rate.

\textbf{Key architectural features:}
\begin{itemize}
    \item \textbf{Dual-track LM}: A ``talker'' module (28-layer Transformer) predicts semantic tokens autoregressively, while a ``code predictor'' generates the 16 RVQ codebook entries in parallel for each frame.
    \item \textbf{Speaker encoder}: Extracts speaker embeddings from reference audio and injects them into the codec embedding layer at a dedicated \texttt{spk\_id} position.
    \item \textbf{Parameter count}: 1.7B (base) and 0.6B (lite).
    \item \textbf{Language support}: 10 languages with Chinese as the primary language.
    \item \textbf{Fine-tuning approach}: Full-parameter SFT (not LoRA), directly learning the speaker embedding.
    \item \textbf{Inference VRAM}: $\sim$20 GB on A100.
    \item \textbf{License}: Apache 2.0 (commercial use permitted).
\end{itemize}

\begin{figure}[htbp]
\centering
\includegraphics[width=0.82\linewidth]{figs/model_architecture_comparison.pdf}
\caption{Architectural comparison between Qwen3-TTS and Fish Audio S2 Pro.}
\label{fig:arch}
\end{figure}

\subsection{Fish Audio S2 Pro (Comparison Model)}

Fish Audio S2 Pro adopts a Dual-AR architecture with a total of $\sim$5B parameters:

\begin{itemize}
    \item \textbf{Slow AR} (4B): Predicts semantic tokens along the time axis.
    \item \textbf{Fast AR} (400M): Fills in residual acoustic details for each frame.
    \item \textbf{Codec}: 10-codebook RVQ at 21 Hz frame rate.
    \item \textbf{Training data}: 10M+ hours across 80+ languages.
    \item \textbf{Zero-shot cloning}: 15--30 seconds of reference audio suffices.
    \item \textbf{Inference VRAM}: $\geq$24 GB.
    \item \textbf{License}: Fish Audio Research License (commercial use requires authorization).
\end{itemize}

\textbf{Why we only evaluate S2 Pro in zero-shot mode:}
\begin{enumerate}
    \item S2 Pro is an RL-trained model; the developers explicitly warn against fine-tuning, as it degrades performance.
    \item By default, fine-tuning only learns speech patterns, not timbre; timbre is captured via the reference prompt at inference time.
    \item The recommended fine-tuning data volume ($\geq$10 hours) far exceeds our available data ($\sim$28 min).
\end{enumerate}

\begin{table}[H]
\caption{Model specification comparison}
\label{tab:models}
\begin{center}
\begin{tabular}{@{}lcc@{}}
\toprule
\textbf{Attribute} & \textbf{Qwen3-TTS 1.7B} & \textbf{Fish S2 Pro} \\
\midrule
Parameters & 1.7B & 5B (4B + 0.4B + codec) \\
Codec & 16-codebook RVQ @ 12Hz & 10-codebook RVQ @ 21Hz \\
Architecture & Transformer Decoder & Dual-AR (Slow + Fast) \\
Training Data & Not disclosed & 10M+ hours \\
Fine-tuning & SFT (full params) & Not recommended \\
Min. Data for FT & 10--30 min & $\geq$10 hours \\
Inference VRAM & $\sim$20 GB & $\geq$24 GB \\
License & Apache 2.0 & Research License \\
\bottomrule
\end{tabular}
\end{center}
\end{table}


%% ============================================================
\section{DATA COLLECTION AND PREPROCESSING}
%% ============================================================

\subsection{Data Sources}

Audio data was collected from Bilibili, an online video platform, targeting a specific character voice actor (``Qinche''). The final training set draws from the following sources:

\begin{table}[htbp]
\caption{Data source summary}
\label{tab:datasource}
\begin{center}
\begin{tabular}{@{}llcp{5.5cm}@{}}
\toprule
\textbf{Source} & \textbf{Type} & \textbf{Samples} & \textbf{Description} \\
\midrule
\texttt{qinche\_pure\_p01$\sim$p11} & Clean single-speaker & 296 & Direct recordings of the target speaker; highest quality \\
\texttt{qinche\_02}                 & Mixed audio           & 77  & Contains other speakers; filtered via Pyannote diarization \\
\texttt{qinche\_call\_p01$\sim$p06} & Call recordings       & 64  & Phone-call style audio; usable after filtering \\
\texttt{qinche\_01}                 & Mixed audio           & 38  & Lower quality; partially retained after strict filtering \\
\bottomrule
\end{tabular}
\end{center}
\end{table}

\subsection{Preprocessing Pipeline}

The preprocessing pipeline consists of six sequential stages, as illustrated in Figure~\ref{fig:pipeline}.

\begin{figure}[htbp]
\centering
\includegraphics[width=0.88\linewidth]{figs/preprocessing_pipeline.pdf}
\caption{End-to-end data preprocessing pipeline. Each stage filters or transforms the data, with quality gates between stages.}
\label{fig:pipeline}
\end{figure}

\begin{enumerate}
    \item \textbf{Source Separation (Demucs)}: Separates vocal tracks from background music and noise using the Demucs neural source separation model.
    \item \textbf{Speaker Diarization (Pyannote)}: For mixed-source audio, identifies and segments speech by individual speakers using Pyannote's speaker diarization pipeline.
    \item \textbf{Voice Activity Detection (Silero-VAD)}: Segments continuous speech into utterance-level clips of 1--15 seconds duration.
    \item \textbf{ASR Transcription (WhisperX)}: Generates Chinese text transcriptions with word-level timestamps.
    \item \textbf{Text Normalization (OpenCC)}: Converts Traditional Chinese characters to Simplified Chinese.
    \item \textbf{Speaker Purification}: Computes speaker embeddings (via Pyannote) and filters segments by cosine similarity to a reference embedding, with random train/test split.
\end{enumerate}

\subsection{Audio Normalization}

After purification, each audio segment undergoes:
\begin{itemize}
    \item \textbf{RMS Loudness Normalization}: Normalizes to a target RMS level for consistent volume across all training samples.
    \item \textbf{Trailing Silence Padding}: Appends 1 second of silence to each segment, as required by the Qwen3-TTS tokenizer for proper end-of-sequence detection.
\end{itemize}

\subsection{Dataset Statistics}

\begin{table}[H]
\caption{Final dataset statistics}
\label{tab:dataset}
\begin{center}
\begin{tabular}{@{}lcccc@{}}
\toprule
\textbf{Split} & \textbf{Samples} & \textbf{Total Duration} & \textbf{Range} & \textbf{Mean Duration} \\
\midrule
Training & 664 & $\sim$28 min & 1.0s -- 13.3s & $\sim$2.5s \\
Test & 18 & $\sim$1.7 min & 4.7s -- 11.7s & 5.8s \\
Reference & 5 clips & --- & 3s -- 10s & --- \\
\bottomrule
\end{tabular}
\end{center}
\end{table}

\begin{figure}[htbp]
\centering
\includegraphics[width=0.62\linewidth]{figs/dataset_duration_distribution.pdf}
\caption{Distribution of audio segment durations in the training set.}
\label{fig:duration}
\end{figure}


%% ============================================================
\section{TRAINING METHODOLOGY}
%% ============================================================

\subsection{Training Pipeline}

The Qwen3-TTS fine-tuning pipeline consists of three sequential stages:
\begin{enumerate}
    \item \textbf{Audio Code Extraction} (\texttt{prepare\_data.py}): Uses the Qwen3-TTS Tokenizer (12Hz) to extract audio codes (16 codebooks) from normalized audio, producing \texttt{train\_with\_codes.jsonl}.
    \item \textbf{SFT Training} (\texttt{sft\_12hz.py}): Full-parameter supervised fine-tuning using HuggingFace Accelerate with bf16 mixed precision and Flash Attention 2.
    \item \textbf{Speaker Embedding Accumulation}: Averages speaker embeddings across all training samples per epoch and saves them alongside the model checkpoint.
\end{enumerate}

\subsection{Hyperparameter Configuration}

\begin{table}[H]
\caption{Training hyperparameters for the best configuration (v5)}
\label{tab:hparams}
\begin{center}
\begin{tabular}{@{}ll@{}}
\toprule
\textbf{Parameter} & \textbf{Value} \\
\midrule
Base Model & Qwen3-TTS-12Hz-1.7B-Base \\
Learning Rate & $2 \times 10^{-6}$ \\
Effective Batch Size & 32 (batch 2 $\times$ gradient accumulation 16) \\
Epochs & 10 \\
Scheduler & Linear warmup (10\%) + Cosine decay (min $0.1 \times$ LR) \\
Precision & bf16 (mixed precision) \\
Attention & Flash Attention 2 \\
Hardware & 1$\times$ NVIDIA A100-SXM4-80GB \\
\bottomrule
\end{tabular}
\end{center}
\end{table}

\subsection{Experiment Iterations}

We conducted five major training iterations (v1--v5), each addressing issues discovered in the previous version.

\begin{figure}[htbp]
\centering
\includegraphics[width=0.88\linewidth]{figs/training_iterations_timeline.pdf}
\caption{Evolution of training iterations showing key changes and their impact on SIM\_gt.}
\label{fig:iterations}
\end{figure}

\subsubsection{v1 --- Initial Attempt (Abandoned)}
Learning rate $2 \times 10^{-5}$ with no scheduler, 387 samples ($\sim$14.7 min). Result: SIM\_gt = 0.04--0.10, WER = 2.9--4.4. \textbf{Complete failure} due to excessive learning rate and lack of learning rate scheduling.

\subsubsection{v2 --- Reduced LR + Scheduler}
Learning rate $2 \times 10^{-6}$ with cosine decay, 461 samples ($\sim$19.7 min). Initial evaluation showed SIM\_gt $\approx$ 0.29 due to a critical \textbf{bf16 precision bug} in speaker embedding accumulation. After fixing (see Section~\ref{sec:bf16bug}), SIM\_gt improved to \textbf{0.6964} at epoch 6 --- a \textbf{+141\%} improvement.

\subsubsection{v3 --- Higher Learning Rate Exploration}
Tested LR = $1 \times 10^{-5}$ (5$\times$ v2). After bf16 fix, SIM\_gt = 0.66--0.68, slightly lower than v2. Conclusion: LR = $2 \times 10^{-6}$ remains optimal.

\subsubsection{v4 --- Post-Bug-Fix Formal Training}
Same configuration as v2 but with the bf16 bug fixed from the start. SIM\_gt = 0.6783 (epoch 7), confirming the fix's correctness.

\subsubsection{v5 --- Data Augmentation (Current Best)}
Expanded dataset from 461 to \textbf{664 samples} ($\sim$28 min, +44\%). SIM\_gt improved to \textbf{0.6892} at epoch 3 with CUDA Graph inference.

\begin{figure}[htbp]
\centering
\includegraphics[width=0.72\linewidth]{figs/sim_gt_across_versions.pdf}
\caption{Best SIM\_gt achieved across training iterations, showing the impact of bug fixes and data augmentation.}
\label{fig:sim_versions}
\end{figure}

\subsection{Critical Bug Fix: bf16 Speaker Embedding Precision}
\label{sec:bf16bug}

A critical bug was discovered in the speaker embedding accumulation process:

\begin{issues}[\textbf{The bf16 Precision Bug}]
During training with \texttt{mixed\_precision="bf16"}, the \texttt{spk\_embed\_sum} tensor was accumulated in bf16 precision. The speaker encoder outputs embeddings with norm $\sim$17, but after accumulating hundreds of samples, the saved embedding norm dropped to only 2.98 due to bf16's limited dynamic range ($\sim$3.4 significant digits).

\textbf{Fix}: Convert the accumulator to fp32 via \texttt{.float()} during accumulation, then cast back to bf16 only at save time. This single fix improved SIM\_gt from 0.29 to 0.70 (\textbf{+141\%}).
\end{issues}

\begin{figure}[htbp]
\centering
\includegraphics[width=0.62\linewidth]{figs/bf16_bug_impact.pdf}
\caption{Impact of the bf16 precision bug on speaker similarity. The fix dramatically improved SIM\_gt from $\sim$0.29 to $\sim$0.70 across all training versions.}
\label{fig:bf16}
\end{figure}


%% ============================================================
\section{EVALUATION}
%% ============================================================

\subsection{Evaluation Metrics}

\begin{table}[H]
\caption{Evaluation metrics definitions}
\label{tab:metrics}
\begin{center}
\begin{tabular}{@{}lp{7.5cm}p{3.8cm}@{}}
\toprule
\textbf{Metric} & \textbf{Definition} & \textbf{Target} \\
\midrule
SIM\_gt & Cosine similarity between generated and ground-truth speaker embeddings (Pyannote) & Higher is better ($>$0.6 = usable) \\
SIM\_ref & Cosine similarity between generated and reference audio speaker embeddings & Higher is better \\
WER & Word Error Rate of ASR transcription vs.\ original text (WhisperX + jiwer) & Lower is better ($<$0.1 = excellent) \\
RTF & Real-Time Factor (generation time / audio duration) & $<$1.0 = real-time \\
\bottomrule
\end{tabular}
\end{center}
\end{table}

\subsection{Comprehensive Results}

\begin{table}[H]
\caption{Full comparison of all evaluated models and configurations}
\label{tab:results}
\begin{center}
\small
\begin{tabular}{@{}lccccc@{}}
\toprule
\textbf{Model / Configuration} & \textbf{SIM\_gt}$\uparrow$ & \textbf{SIM\_ref}$\uparrow$ & \textbf{WER}$\downarrow$ & \textbf{RTF}$\downarrow$ & \textbf{Inference} \\
\midrule
Qwen3 1.7B Zero-shot & \textbf{0.7104} & 0.7563 & 0.1203 & 2.401 & Native \\
\midrule
Qwen3 SFT v5 ep3 & \textbf{0.6892} & 0.7161 & 0.0425 & \textbf{0.357} & CUDA Graph \\
Qwen3 SFT v5 ep3 & 0.6885 & \textbf{0.7281} & 0.0451 & 2.542 & Native \\
Qwen3 SFT v5 ep5 & 0.6791 & 0.6998 & \textbf{0.0279} & 0.356 & CUDA Graph \\
Qwen3 SFT v5 ep9 & 0.6795 & 0.7180 & 0.0269 & 0.354 & CUDA Graph \\
\midrule
Qwen3 SFT v2-fixed ep6 & 0.6964 & \textbf{0.7653} & 0.0365 & 2.579 & Native \\
Qwen3 SFT v4 ep7 & 0.6783 & 0.7034 & 0.0530 & 0.351 & CUDA Graph \\
\midrule
Fish S2 Pro ZS (compiled) & 0.6627 & 0.6824 & \textbf{0.0209} & 0.639 & torch.compile \\
Fish S2 Pro ZS & 0.6647 & 0.6910 & 0.0244 & 4.315 & Native \\
\bottomrule
\end{tabular}
\end{center}
\end{table}

\begin{figure}[htbp]
\centering
\includegraphics[width=0.72\linewidth]{figs/radar_comparison.pdf}
\caption{Radar chart comparing key configurations across four metrics. Metrics are normalized so that larger area indicates better overall performance.}
\label{fig:radar}
\end{figure}


\subsection{Per-Sample Analysis}

\begin{figure}[htbp]
\centering
\includegraphics[width=0.88\linewidth]{figs/per_sample_sim_gt.pdf}
\caption{Per-sample SIM\_gt comparison across key configurations (18 test samples). Most samples show consistent improvement from fine-tuning, with some variation on challenging samples.}
\label{fig:persample}
\end{figure}

\begin{figure}[htbp]
\centering
\includegraphics[width=0.88\linewidth]{figs/per_sample_wer.pdf}
\caption{Per-sample WER comparison. Fine-tuned models achieve near-zero WER on most samples, with occasional outliers on texts containing rare characters or mixed Traditional/Simplified Chinese.}
\label{fig:persamplewer}
\end{figure}

\subsection{Key Findings}

\begin{issues}[\textbf{Key Findings from Evaluation}]
\begin{enumerate}
    \item \textbf{Data augmentation is effective}: v5 (664 samples, 28 min) vs. v4 (461 samples, 19.7 min): SIM\_gt improved from 0.6783 to 0.6892 (+0.011).
    \item \textbf{bf16 precision bug was the root cause of low SIM}: Fix improved SIM\_gt from 0.29 $\to$ 0.70 (+141\%).
    \item \textbf{CUDA Graph inference is lossless}: v5 ep3 native SIM\_gt = 0.6885 $\approx$ CUDA Graph SIM\_gt = 0.6892.
    \item \textbf{LR = $2\times10^{-6}$ is optimal}: v3 (LR = $1\times10^{-5}$) showed no improvement over v2.
    \item \textbf{Fish S2 Pro excels in WER}: Best WER = 0.0209, likely due to its massive 10M+ hour training corpus.
    \item \textbf{Fine-tuned Qwen3 dominates in SIM}: SIM\_gt = 0.689 vs Fish S2 Pro's 0.665 for the target speaker.
    \item \textbf{Early stopping at epoch 3 is optimal}: Later epochs show diminishing returns or degradation.
\end{enumerate}
\end{issues}


%% ============================================================
\section{INFERENCE ACCELERATION}
%% ============================================================

\subsection{Verified Acceleration Methods}

\begin{table}[H]
\caption{Inference acceleration benchmark results}
\label{tab:accel}
\begin{center}
\begin{tabular}{@{}lccccc@{}}
\toprule
\textbf{Method} & \textbf{Model} & \textbf{Original RTF} & \textbf{Optimized RTF} & \textbf{Speedup} & \textbf{SIM Loss} \\
\midrule
FasterQwen3TTS (CUDA Graph) & Qwen3-TTS & 2.542 & \textbf{0.357} & \textbf{7.1$\times$} & None \\
Fish S2 Pro \texttt{--compile} & Fish S2 Pro & 4.315 & \textbf{0.639} & \textbf{6.75$\times$} & None \\
Flash Attention 2 & Both & --- & --- & --- & Baseline \\
Streaming + CUDA Graph & Qwen3-TTS & 2.542 & 0.395 & 5.88$\times$ & None \\
\bottomrule
\end{tabular}
\end{center}
\end{table}

\begin{figure}[htbp]
\centering
\includegraphics[width=0.72\linewidth]{figs/rtf_comparison_bar.pdf}
\caption{RTF comparison across all models and acceleration methods. The dashed line at RTF = 1.0 marks the real-time threshold.}
\label{fig:rtf}
\end{figure}

\subsubsection{FasterQwen3TTS (Recommended)}

The \texttt{faster-qwen3-tts} library uses CUDA Graph capture to eliminate kernel launch overhead:
\begin{itemize}
    \item Captures separate CUDA Graphs for the talker and code predictor modules during a warmup phase.
    \item Achieves \textbf{RTF = 0.357} on A100 (7.1$\times$ speedup).
    \item Streaming mode: RTF = 0.395, TTFA (Time to First Audio) = 305ms.
    \item Directly compatible with fine-tuned checkpoints via \texttt{from\_pretrained}.
    \item No dependency on vLLM, Triton, or custom CUDA kernels.
\end{itemize}

\subsubsection{Fish S2 Pro \texttt{--compile}}

Fish Speech natively supports \texttt{torch.compile} for the \texttt{decode\_one\_token\_ar} function:
\begin{itemize}
    \item RTF: 4.315 $\to$ 0.639 (\textbf{6.75$\times$ speedup}).
    \item First compilation overhead: $\sim$3.5 seconds; subsequent inference stable at $\sim$22 tokens/sec.
    \item GPU memory: 20.1 GB.
\end{itemize}

\subsection{Rejected Acceleration Methods}

\begin{table}[H]
\caption{Acceleration methods that were attempted but rejected}
\label{tab:rejected}
\begin{center}
\small
\begin{tabular}{@{}lp{4cm}p{5.5cm}@{}}
\toprule
\textbf{Method} & \textbf{Result} & \textbf{Rejection Reason} \\
\midrule
INT8 Quantization (bitsandbytes) & SIM\_gt = 0.07, WER = 1.0, RTF = 7.0 & Incompatible with Qwen3-TTS's dual-track LM architecture; quantization completely distorts output \\
Speculative Decoding (0.6B draft) & Acceptance rate: 0--5\% & Un-fine-tuned 0.6B Base token distribution diverges drastically from SFT 1.7B model \\
Self-Speculative (early-exit) & Inter-layer agreement $<$ 7\% & All 28 talker layers perform substantive feature transformation; no redundant layers exist \\
\bottomrule
\end{tabular}
\end{center}
\end{table}

\begin{figure}[htbp]
\centering
\includegraphics[width=0.72\linewidth]{figs/acceleration_summary.pdf}
\caption{Summary of all attempted acceleration strategies. Green indicates successful methods; red indicates rejected methods.}
\label{fig:accel_summary}
\end{figure}

\subsection{Acceleration Conclusions}

CUDA Graph capture reduces RTF from 2.54 to 0.36 (\textbf{7.1$\times$ speedup}), well below real-time, with \textbf{zero quality loss}. Further acceleration attempts (quantization, speculative decoding) all fail due to Qwen3-TTS's unique dual-track language model architecture. The current CUDA Graph solution represents a \textbf{near-optimal acceleration} for this architecture.


%% ============================================================
\section{EPOCH-LEVEL ANALYSIS}
%% ============================================================

\needspace{6\baselineskip}
\subsection{Training Convergence across Epochs}

\begin{figure}[htbp]
\centering
\includegraphics[width=0.60\linewidth]{figs/epoch_sim_wer_v5.pdf}
\caption{SIM\_gt and WER across epochs for v5 (FasterQwen3TTS inference). Epoch 3 achieves the best SIM\_gt while maintaining acceptable WER.}
\label{fig:epoch_v5}
\end{figure}

\needspace{6\baselineskip}
\subsection{Native vs. CUDA Graph Inference Comparison}

\begin{figure}[htbp]
\centering
\begin{subfigure}[b]{0.48\linewidth}
  \centering
  \includegraphics[width=\linewidth]{figs/epoch_comparison_v4_v5.pdf}
  \caption{v4 vs. v5 SIM\_gt per epoch (CUDA Graph).}
  \label{fig:epoch_v4v5}
\end{subfigure}
\hfill
\begin{subfigure}[b]{0.48\linewidth}
  \centering
  \includegraphics[width=\linewidth]{figs/native_vs_fast_v4.pdf}
  \caption{Native vs. CUDA Graph SIM\_gt for v4.}
  \label{fig:native_vs_fast}
\end{subfigure}
\caption{Left: data augmentation benefit (v4 vs. v5) across epochs. Right: CUDA Graph inference is quality-preserving (difference $<$0.02).}
\label{fig:epoch_combined}
\end{figure}


%% ============================================================
\section{CONCLUSIONS AND FUTURE WORK}
%% ============================================================

\subsection{Summary}

This report demonstrates a complete, reproducible pipeline for Chinese TTS voice cloning:

\begin{enumerate}
    \item \textbf{Data}: 664 curated samples ($\sim$28 min) from online sources, processed through a 6-stage pipeline (Demucs $\to$ Pyannote $\to$ Silero-VAD $\to$ WhisperX $\to$ OpenCC $\to$ speaker purification).
    \item \textbf{Training}: Qwen3-TTS-1.7B SFT with LR = $2\times10^{-6}$, cosine scheduler, bf16 mixed precision. Best at epoch 3.
    \item \textbf{Quality}: SIM\_gt = 0.689, WER = 4.25\%, competitive with the 5B-parameter Fish S2 Pro in zero-shot mode.
    \item \textbf{Speed}: RTF = 0.357 (7.1$\times$ real-time) via FasterQwen3TTS CUDA Graph capture, with negligible quality loss.
\end{enumerate}

\subsection{Production Recommendation}

\begin{issues}[\textbf{Deployment Recommendations}]
\begin{itemize}
    \item[-] \textbf{Production}: Qwen3 SFT v5 ep3 + FasterQwen3TTS (RTF = 0.36, SIM\_gt = 0.69)
    \item[-] \textbf{Maximum quality}: Qwen3 SFT v5 ep3 native inference (SIM\_gt = 0.69, RTF = 2.54)
    \item[-] \textbf{Zero-shot alternative}: Fish S2 Pro + compile (RTF = 0.64, no fine-tuning required)
\end{itemize}
\end{issues}

\subsection{Future Work}

\begin{enumerate}
    \item \textbf{Data expansion}: Current 28 min is at the lower bound of recommended range (30--60 min). Additional high-quality data could further improve SIM.
    \item \textbf{Reference audio optimization}: Systematic study of how different reference audio combinations affect synthesis quality.
    \item \textbf{Subjective evaluation}: Conduct Mean Opinion Score (MOS) tests focusing on timbre similarity, naturalness, prosody, and emotional expression.
    \item \textbf{Service deployment}: Leverage FasterQwen3TTS's built-in server mode or vLLM-Omni for high-concurrency API serving.
\end{enumerate}


\end{document}
